\documentclass[11pt]{amsart}
\usepackage[OT1]{fontenc}
\usepackage[margin=1.08in]{geometry}
\usepackage{amsmath,amssymb,amsthm,mathtools}
\usepackage{enumitem}
\usepackage{microtype}
\usepackage{tikz-cd}
\usepackage{xcolor}
\definecolor{linkblue}{RGB}{28,65,108}
\usepackage[colorlinks=true,linkcolor=linkblue,citecolor=linkblue,urlcolor=linkblue]{hyperref}
\numberwithin{equation}{section}
\newtheorem{theorem}{Theorem}[section]
\newtheorem{proposition}[theorem]{Proposition}
\newtheorem{lemma}[theorem]{Lemma}
\newtheorem{corollary}[theorem]{Corollary}
\theoremstyle{definition}
\newtheorem{definition}[theorem]{Definition}
\theoremstyle{remark}
\newtheorem{remark}[theorem]{Remark}
\setlist[enumerate]{label=\textup{(\roman*)},leftmargin=2.2em}
\setlist[itemize]{leftmargin=1.6em}
\allowdisplaybreaks[1]
\title{Minimal models in numerical dimension one}
\author{OpenAI}
\date{September 24, 2026}
\hypersetup{
  pdftitle={Minimal models in numerical dimension one},
  pdfauthor={OpenAI},
  pdfsubject={Minimal models, numerical dimension, and surface intersection theory}
}
\begin{document}
\begin{abstract}
We resolve the numerical-dimension-one case of the minimal-model conjecture
for smooth connected complex projective varieties of dimension at least three.
If $K_X$ is pseudo-effective and $\kappa_\sigma(X,K_X)=1$, with
$\kappa_\sigma$ defined by section growth with a fixed ample twist, then $X$
admits a projective $\mathbb Q$-factorial terminal minimal model.
\end{abstract}
\maketitle
\section{Introduction}

The minimal-model problem asks whether a smooth projective variety whose
canonical divisor is pseudo-effective admits a birational model with nef
canonical divisor. Here pseudo-effective means that the numerical class lies
in the closure of the cone generated by effective divisors, and nef means
nonnegative degree on every integral curve. The models allowed by the minimal
model program are normal and may have terminal singularities. A minimal
model is required to improve canonical discrepancies and to extract no
divisors, as well as to have nef canonical divisor. These requirements retain
the birational information of the original variety; the precise comparison
used here is part of Theorem~\ref{main:theorem}.

We prove existence in the case where the canonical divisor has numerical
dimension one, using its growth of sections with a fixed ample twist. The
argument begins with minimal model programs for positive boundary
perturbations, but its final step is a criterion on surfaces. In particular,
the proof does not require a general minimal-model existence theorem.

\subsection{The invariant and the main theorem}

For a Cartier divisor \(D\) on a smooth projective \(n\)-fold \(X\),
we use Nakayama's section-growth invariant \cite{Nak04} in the following form:
\begin{equation}\label{main:sigma}
\kappa_\sigma(X,D)=
\max\left\{k\in\{0,\ldots,n\}:
\begin{array}{l}
\text{there is an ample Cartier divisor \(A\) with}\\[2pt]
\displaystyle\limsup_{m\to\infty}
\frac{h^0(X,\mathcal O_X(mD+A))}{m^k}>0
\end{array}\right\}.
\end{equation}
The maximum is \(-\infty\) if the set is empty, and \(m\) ranges over
positive integers. The twist \(A\) is fixed before taking the limsup. Thus
\(\kappa_\sigma(X,D)=1\) rules out positive quadratic limsup for
\emph{every} fixed ample Cartier twist. We will use exactly this consequence,
including when quadratic growth is obtained only along the multiples of one
fixed positive integer. It does not assert an upper bound of the form
\(h^0(X,mD+A)=O(m)\).

This growth convention should be distinguished from the intersection
formula for a nef divisor,
\[
\nu(D)=\max\{k\in\{0,\ldots,n\}:D^k\cdot H^{n-k}>0\},
\]
where \(H\) is ample. We do not use that formula for \(K_X\) before
constructing a nef model, nor do we need a comparison theorem between
numerical-dimension conventions.

\begin{samepage}
\begin{theorem}\label{main:theorem}
Let \(X\) be a smooth connected complex projective variety of dimension
\(n\geq3\). Suppose that \(K_X\) is pseudo-effective and
\(\kappa_\sigma(X,K_X)=1\), with the convention
\eqref{main:sigma}. Then there are a normal projective
\(\mathbb Q\)-factorial terminal variety \(Y\) with \(K_Y\) nef and
a finite sequence
\(\phi:X\dashrightarrow Y\) of \(K\)-negative divisorial contractions
and flips. The inverse of \(\phi\) contracts no prime divisor. On a
common smooth projective resolution \(p:W\to X\), \(q:W\to Y\),
compatible canonical divisors satisfy
\begin{equation}\label{main:comparison}
p^*K_X=q^*K_Y+E,
\end{equation}
where \(E\) is an effective \(q\)-exceptional \(\mathbb Q\)-divisor
whose support contains the strict transform of every prime divisor of
\(X\) contracted by \(\phi\).
\end{theorem}
\end{samepage}

The theorem gives a positive resolution of the numerical-dimension-one
case of the minimal-model conjecture for smooth complex projective varieties
of dimension at least three. In particular it applies to smooth fivefolds,
with no hypothesis on \(\chi(X,\mathcal O_X)\).

A \emph{good} minimal model additionally requires a positive multiple of
\(K_Y\) to be generated by global sections. This is the semiampleness
conclusion of abundance. Nonvanishing asks for a nonzero section of a
positive multiple of the canonical divisor. Neither conclusion is asserted by
Theorem~\ref{main:theorem} alone; the fixed ample twist in
\eqref{main:sigma} is retained throughout the section-growth argument.
The following separate consequence uses log abundance.

\begin{corollary}[Good minimal model]\label{cor:good-model}
Under the hypotheses of Theorem~\ref{main:theorem}, the same
\(\mathbb Q\)-factorial terminal endpoint \(Y\) has semiample \(K_Y\);
hence \(Y\) is a good minimal model of \(X\).
\end{corollary}
\begin{proof}
Apply the log abundance theorem \cite[Theorem~1.1]{OpenAILogAbundance2026}
to the projective lc pair \((Y,0)\) over \(\mathbb C\), whose
\(\mathbb Q\)-Cartier canonical divisor is nef.
\end{proof}

\subsection{The problem and the methods in the literature}

The three-dimensional program developed through work of Reid, Benveniste,
Kawamata, Shokurov, Koll\'ar, Miyaoka and others, with Mori's flip theorem
providing a central step in minimal-model existence \cite{Mori88};
see \cite{KM98} for the broader development.
The minimal model program replaces a variety by successive contractions
and flips along canonical-negative extremal rays; the standard foundations,
including the singularities and discrepancy comparisons used here, are
presented in Koll\'ar--Mori \cite{KM98}. Birkar's formulation for log pairs
\cite[Conjecture~1.1]{Birkar10} places existence of log minimal models
alongside the Mori fiber space alternative. A decisive advance was the work
of Birkar, Cascini, Hacon and McKernan \cite{BCHM10}, which established
minimal-model existence for varieties of general type and finite minimal
model programs with scaling for klt pairs with big boundary. The positive
perturbations in our proof lie within precisely this established range.
Their boundary transforms remain big; they need not remain ample.

At the other end of the numerical spectrum, Nakayama developed the
section-growth invariants and divisorial Zariski decompositions that underlie
numerical approaches to the canonical divisor \cite{Nak04}. Druel proved
termination results for directed programs in numerical dimension zero
\cite[Corollary~3.4]{Druel11}, and Gongyo extended this method to prove
minimal-model existence for \(\mathbb Q\)-factorial dlt pairs of numerical
log Kodaira dimension zero
\cite[Theorem~1.1]{Gongyo11}. The problem here is to pass
from arbitrarily small positive perturbations to the canonical divisor in
the next numerical dimension. A finite program for each positive parameter
does not by itself terminate their concatenation at parameter zero.

Our first reduction also has a specific methodological predecessor.
Alexeev, Hacon and Kawamata use fourth homology to control certain
four-dimensional flips \cite[Lemma~3.1 and its proof]{AHK07}. We use the
same kind of decrease after a discrepancy argument has removed
codimension-two components on the flipped side. In arbitrary dimension
\(n\), the relevant finite-dimensional space is the span of algebraic
\((n-2)\)-cycle classes in degree \(2n-4\). The proof of
Lemma~\ref{mmp:stabilization} gives the needed statement directly; the
four-dimensional result is a precedent for the method, not an imported
termination theorem in all dimensions.

The second reduction combines two established strands of positivity.
Nadel's multiplier-ideal vanishing theorem \cite{Nadel90} supplies the
vanishing used to extend sections from a very general surface while allowing the perturbation to shrink with the section degree;
we use the formulation in Lazarsfeld
\cite[Theorem~9.4.8]{Laz04II}, \cite[Theorem~2.4]{LazMI}.
Surface Riemann--Roch then converts positive square into quadratic section
growth. A different surface, through a point of a hypothetical negative
curve, turns a vanishing-order estimate into a fixed exceptional divisor.
The Hodge index theorem makes its square uniformly negative. The surface
intersection calculations are classical \cite[Chapter~V, Section~1]{Hartshorne77};
our use of them keeps one surface resolution and one exceptional curve fixed
while the rational perturbations vary. This uniformity is what produces a
contradiction.

For comparison with results concerning sections of the unperturbed
canonical divisor, Lazi\'c and Peternell prove nonvanishing for an
\emph{already minimal} projective terminal variety with numerical dimension
one and nonzero Euler characteristic \cite[Theorem~6.7]{LP18}.
Liu and Xu prove existence of a good minimal model for smooth projective
varieties of dimension at most five and numerical dimension at most one,
assuming nonnegative Kodaira dimension \cite[Theorem~1.2]{LiuXu25}.
They also obtain a good minimal model for a smooth projective variety of
dimension at most four with nonzero Euler characteristic and
\(0\leq\kappa_\sigma(X,K_X)\leq1\)
\cite[Corollary~5.2]{LiuXu25}.
These hypotheses and conclusions differ from the existence statement above.
We do not use their nonvanishing or abundance results as proof inputs.

\subsection{The route through the proof}

The first step is to study a sequence of shrinking perturbations on one
fixed model.
Starting with an effective ample rational boundary \(B\), we run finite
minimal model programs for decreasing positive coefficients of its
transform. Every step is \(K\)-negative. Discrepancy and cycle-class
arguments show that only finitely many steps change a locus of codimension
at most two; they do not assert termination of the concatenated programs.
Proposition~\ref{mmp:approximation} therefore supplies a finite canonical
prefix \(X\dashrightarrow Y\) such that multiples of
\[
M_j=K_Y+t_jB_Y,\qquad t_j>0,\quad t_j\longrightarrow0,
\]
are generated outside closed subsets \(Z_j\subset Y\) of codimension
at least three. Here \(B_Y\) is the transform of \(B\); the generated
multiple and the bad set may depend on \(j\).

Two surfaces then show that \(K_Y\) is nef. Fix a very ample divisor
\(H\) on \(Y\). A very general complete intersection surface \(S\)
avoids every \(Z_j\) and the singular locus. The restrictions
\(M_j|_S\) are semiample, so their limit \(K_Y|_S\) is nef.
Lemma~\ref{num:growth} uses multiplier-ideal vanishing to extend sections
from this fixed surface with one fixed twist. Positive square
\(K_Y^2\cdot H^{n-2}\) would consequently give quadratic section growth.
Section~\ref{main:proof-section} transfers that growth to the original
\(X\), where \(\kappa_\sigma(X,K_X)=1\) rules it out. Thus the square
is zero.

This first surface need not meet a curve of negative \(K_Y\)-degree.
To exclude such a curve, including one in the singular locus, take another
complete intersection from \(|H|\) through one of its points, and choose
a surface component through that point. On every component the restriction
of \(M_j\) has nonnegative square; these squares, counted with
multiplicity, sum to \(M_j^2\cdot H^{n-2}\to0\). Thus the square on
the chosen surface also tends to zero.
For all sufficiently small perturbations, an ambient jet estimate on a
resolution of \(Y\) forces sections of Cartier multiples \(kM_j\) to vanish to order at
least \(ck\) at one fixed point above the curve, with \(c>0\) independent
of \(j\) and \(k\). On a fixed resolution of the second surface, this
produces an exceptional fixed part. The Hodge index theorem bounds its
square above by \(-\epsilon k^2\), with \(\epsilon>0\) fixed, while
the residual moving system has nonnegative square. Thus the perturbations
have uniformly positive square on that surface, contradicting its
zero-square limit. This is the nefness
criterion of Proposition~\ref{num:nef}. The finite prefix already supplies
the remaining birational and canonical-comparison assertions of
Theorem~\ref{main:theorem}.

\subsection{Conventions}

All varieties and morphisms are over \(\mathbb C\), and all birational
models in the proof are projective. We use compatible canonical divisors.
Log discrepancies are normalized by
\[
a(F;U,\Delta)=1+\operatorname{coeff}_F
\bigl(K_{\widetilde U}-r^*(K_U+\Delta)\bigr),
\]
where \(r:\widetilde U\to U\) is a resolution carrying the prime
divisor \(F\). Terminality at zero boundary means
\(a(F;U,0)>1\) for every exceptional prime \(F\) over \(U\);
an original prime divisor has log discrepancy \(1\).
For an effective rational boundary, klt means that all log discrepancies
are positive \cite[Chapter~2]{KM98}.

A birational map \emph{extracts no divisors} if its inverse contracts no
prime divisor. Intersection numbers of rational Cartier divisors are defined
after clearing denominators. On an integral surface that need not be normal,
we use Cartier intersections with its fundamental cycle
\cite[Chapter~2]{Fulton98}. Characteristic-zero projective resolutions and
common resolutions are used throughout \cite{Hironaka64}.

\section{Positive perturbations on a fixed model}
\label{mmp:section}

The construction in this section does not use the numerical-dimension
hypothesis.  It provides one model on which small positive perturbations
of the canonical divisor have base locus of codimension at least three.

\begin{proposition}\label{mmp:approximation}
Let $X$ be a smooth projective variety of dimension $n\geq 3$ with
pseudo-effective $K_X$.  There are an effective ample $\mathbb Q$-divisor
$B$ on $X$, an integer $s\geq 0$, and a finite sequence of
$K$-negative divisorial contractions and $K$-flips
\[
 X\dashrightarrow Y
\]
with the following properties.
\begin{enumerate}[label=\textup{(\roman*)}]
\item $Y$ is projective, $\mathbb Q$-factorial, and terminal.  The map
$X\dashrightarrow Y$ extracts no divisors.
\item Put $L=K_Y$, let $B_Y$ be the transform of $B$, and set
\[
 t_j=2^{-j},\qquad M_j=L+t_jB_Y\quad (j>s).
\]
For each $j>s$ there is a closed subset $Z_j\subset Y$ of codimension
at least three such that $|kM_j|$ is defined by a Cartier divisor and
is base point free on $Y\setminus Z_j$ for every sufficiently divisible
positive integer $k$.
\item On a common smooth projective resolution
$p:W\to X$, $q:W\to Y$, compatible canonical divisors satisfy
\begin{equation}\label{mmp:comparison}
 p^*K_X=q^*K_Y+E,
\end{equation}
where $E\geq 0$ is $q$-exceptional and its support contains the strict
transform of every prime divisor of $X$ contracted by $X\dashrightarrow Y$.
\end{enumerate}
\end{proposition}

Here and below, ``sufficiently divisible'' allows the required divisor
index to depend on $j$.  No uniform Cartier index for an infinite
sequence of models will be needed.

\subsection{Discrepancies and codimension two}

Write $a(F,U)=a(F;U,0)$ with the normalization fixed in the
introduction.  We record the strictness statement needed for the
finiteness argument.

\begin{lemma}\label{mmp:strict}
Suppose that $U\dashrightarrow U^+$ is a $K$-negative divisorial
contraction or a $K$-flip between normal projective
$\mathbb Q$-Gorenstein varieties.  Write $h:U\to R$ for the
contraction and $h^+:U^+\to R$ for the other morphism, with $h^+$
the identity in the divisorial case.  For every prime divisor $F$
over these varieties,
\[
 a(F,U)\leq a(F,U^+).
\]
The inequality is strict if the center of $F$ on $R$ is contained
in the locus where either $h$ or $h^+$ has a positive-dimensional
fiber.
\end{lemma}

\begin{proof}
Take a common smooth projective resolution carrying $F$:
\[
\begin{tikzcd}[column sep=large,row sep=small]
 & W \arrow[dl,"p"'] \arrow[dr,"q"] & \\
 U \arrow[dr,"h"'] \arrow[rr,dashed] & & U^+ \arrow[dl,"h^+"] \\
 & R. &
\end{tikzcd}
\]
Write $g=h\circ p=h^+\circ q:W\to R$ and set
$D=p^*K_U-q^*K_{U^+}$.  This divisor is $q$-exceptional,
since the step extracts no divisors.  On every curve contracted by
$g$, the divisors $-p^*K_U$ and $q^*K_{U^+}$ have nonnegative
degree; in the divisorial case the second has degree zero.
Thus $-D$ is $g$-nef, and hence $q$-nef.  The negativity lemma
gives $D\geq 0$; see \cite[Lemmas~3.38--3.39]{KM98}.

To verify the support assertion, let $r\in R$ be
a closed point with a positive-dimensional fiber on either side.
There is a curve in $g^{-1}(r)$ mapping onto a curve in that fiber:
take a component dominating that curve and cut it by sufficiently
general ample divisors.  The relative ampleness signs give $D\cdot C<0$
for this lifted curve.  Consequently $g^{-1}(r)$ meets
$\operatorname{Supp}D$.

We use the fiber-support form of the negativity lemma
\cite[Lemma~3.39]{KM98}: the support of an effective divisor
antinef over a proper birational morphism to a normal variety either
contains or is disjoint from each fiber. In the present setting it can
also be seen directly on the reduced irreducible components of the
fiber. If such a component $T$ is not contained in $\operatorname{Supp}D$
but meets it, a Cartier multiple of $D$ restricts to a nonzero effective
Cartier divisor on the integral projective variety $T$. Its degree against
sufficiently many ample hyperplanes is positive, producing a curve
contracted by $g$ with positive $D$-degree. This contradicts antinefness.
This argument does not require $T$ to be normal or the fiber to be reduced:
only its support is at issue. Components of dimension zero already lie
in the support if they meet it. The fibers of $g$ are connected because
$R$ is normal and $g$ is proper and birational
\cite[Corollary~III.11.4]{Hartshorne77}. Containment therefore propagates
through their reduced irreducible components.

It follows that every fiber in question is contained in
$\operatorname{Supp}D$.  If the center of $F$ on $R$ is contained
in this locus, then $F\subset\operatorname{Supp}D$ on our chosen
resolution.  Finally,
\[
 a(F,U^+)-a(F,U)=\operatorname{coeff}_F D,
\]
which proves both assertions.
\end{proof}

\begin{lemma}[Stabilization in codimension two]\label{mmp:stabilization}
Let
\[
 X=U_0\dashrightarrow U_1\dashrightarrow U_2\dashrightarrow\cdots
\]
be a finite or infinite sequence of $K$-negative divisorial
contractions and $K$-flips of projective $\mathbb Q$-factorial
varieties, starting from a smooth $n$-fold, $n\geq 3$.
After finitely many steps, each remaining map is an isomorphism
outside closed subsets of codimension at least three on both sides.
\end{lemma}

\begin{proof}
The usual MMP properties give preservation of $\mathbb Q$-factoriality,
invariance of the Picard number under a flip, and a drop of one under
an elementary divisorial contraction
\cite[Propositions~3.36--3.37]{KM98}.
There are therefore only finitely many divisorial contractions.
Since the steps extract no divisors, only finitely many prime divisors
of $X$ can disappear; denote this finite set by $\mathcal D$.

All models are terminal.  This follows from
\cite[Corollaries~3.42--3.43]{KM98}, or directly from
Lemma~\ref{mmp:strict}: previously exceptional valuations retain
log discrepancy greater than one, and a newly contracted divisor
increases strictly from log discrepancy one.  Moreover, every
valuation exceptional over the smooth $X$ has integral log
discrepancy at least two.  Discrepancies never decrease along the
sequence.

\smallskip\noindent
\emph{The flipped loci.}
Consider a flip whose target has a codimension-two component $T$
of its flipped locus.  A terminal variety is smooth in codimension
two \cite[Corollary~5.18]{KM98}.  Blowing up $T$ at its generic
smooth point therefore defines a prime valuation $F$ with
\[
 a(F,U_{i+1})=2.
\]
Its center on the contraction base is contained in the locus of
positive-dimensional fibers of the flipped contraction.
Lemma~\ref{mmp:strict} implies
\begin{equation}\label{mmp:witness}
 a(F,X)\leq a(F,U_i)<2=a(F,U_{i+1}).
\end{equation}
Thus $F$ is not exceptional over $X$: it is an original prime divisor.
It is exceptional over $U_{i+1}$, so it belongs to $\mathcal D$.
After this occurrence its discrepancy is at least two forever,
and it cannot satisfy \eqref{mmp:witness} at a later step.
Each flip with such a target component therefore uses a distinct
member of a finite set.  After a finite prefix, all steps are flips
and every flipped locus has codimension at least three.

\smallskip\noindent
\emph{The flipping loci.}
We have removed codimension-two components on the target side. It remains
to remove them on the source side; only then can a general surface be
transported unchanged through every later finite stage. The argument
adapts the homological decrease used in the proof of
\cite[Lemma~3.1]{AHK07}, keeping only classes of algebraic cycles.
Put $d=n-2$ and
define the finite-dimensional real vector space
\[
 \mathcal A_d(V)
 =\operatorname{span}_{\mathbb R}
 \{[T]:T\subset V\text{ is a closed integral $d$-fold}\}
 \subset H_{2d}(V,\mathbb R).
\]
Cycle classes,
Borel--Moore localization, and their compatibility with restriction
to an open subset are as in \cite[Section~19.1]{Fulton98}.
The homology of the compact complex algebraic varieties here is
finite-dimensional and vanishes above their real dimension.

For a remaining flip $U\dashrightarrow U^+$ over $R$, remove the
images in $R$ of both exceptional loci and take inverse images.
This gives a common open subset $\mathcal U$.  Indeed, a proper
birational morphism to a normal variety is an isomorphism over its
quasi-finite locus.  Each exceptional locus consists of
positive-dimensional fibers; its image has dimension at most one
less than its own.  Since the source contraction is small and the
target exceptional locus has dimension at most $n-3$, the closed
complements $Z\subset U$ and $Z^+\subset U^+$ satisfy
\[
 \dim Z\leq d,\qquad \dim Z^+\leq d-1.
\]
The additional points removed outside the exceptional loci lie in
an isomorphism locus and have dimension at most $n-3$.

Restriction gives maps
\[
 \mathcal A_d(U)\longrightarrow
 H^{\mathrm{BM}}_{2d}(\mathcal U,\mathbb R)
 \longleftarrow\mathcal A_d(U^+)
\]
with the same image: intersect an integral $d$-fold with
$\mathcal U$ and take its closure on the other model, while a
$d$-fold contained in the complement restricts to zero.
The right-hand map is injective, by the localization sequence and
$H_{2d}(Z^+,\mathbb R)=0$.
Hence
\begin{equation}\label{mmp:rank}
 \dim\mathcal A_d(U^+)\leq\dim\mathcal A_d(U).
\end{equation}
If the flipping locus has a $d$-dimensional component $T$, its
class belongs to the kernel of the left-hand restriction and is
nonzero: for an ample Cartier divisor $A$,
\[
 \langle c_1(A)^d,[T]\rangle=A^d\cdot T>0.
\]
Then \eqref{mmp:rank} is strict.  A nonnegative integer can drop only
finitely many times, so eventually the flipping loci also have
codimension at least three.  The same construction of the common
open then gives this codimension bound on both complements.
\end{proof}

\subsection{Construction of the perturbations}

\begin{proof}[Proof of Proposition~\ref{mmp:approximation}]
Choose an effective ample rational divisor $B$ for which $(X,B)$
is klt and $K_X+B$ is ample.  For example, take a sufficiently
positive general smooth very ample divisor and multiply it by a
rational number strictly between zero and one.  Since $K_X$ is
pseudo-effective, $K_X+tB$ is big for every rational $t>0$.

We shall repeatedly use a simple consequence of normality.  If a
birational map $X\dashrightarrow U$ extracts no divisors, every
prime divisor of $U$ corresponds to a prime divisor of $X$.
For compatible transformed divisors, the divisorial criterion for
regularity of a rational section consequently gives
\begin{equation}\label{mmp:sections}
 H^0(X,kD)\hookrightarrow H^0(U,kD_U)
\end{equation}
whenever the multiples are Cartier.  Thus effectiveness and bigness
of $B_U$, and bigness of $K_U+tB_U$, persist on every model
constructed below.  In particular, we do not need $B_U$ to be ample.

\smallskip\noindent
\emph{Finite stages.}
Set $t_j=2^{-j}$ for every integer $j\geq0$.
Starting with $V_0=X$, construct $V_{j+1}$ from $V_j$ by running
the MMP for
\[
 K_{V_j}+\Delta_j,\qquad
 \Delta_j=t_{j+1}B_{V_j},
\]
with scaling of
$C_j=(t_j-t_{j+1})B_{V_j}$.
The precise input is \cite[Corollary~1.4.2]{BCHM10}:
for a projective $\mathbb Q$-factorial klt pair $(V,\Delta)$
with big boundary $\Delta$, and $C\geq0$ such that
$(V,\Delta+C)$ is klt and $K_V+\Delta+C$ is nef,
the MMP with scaling of $C$ terminates.
Here the inductive hypothesis is that $(V_j,t_jB_{V_j})$ is
klt and its adjoint is nef.  The smaller pair is klt because
$B_{V_j}$ is effective; its boundary is big by
\eqref{mmp:sections}; and the larger pair is exactly the one in
the inductive hypothesis.

The adjoint remains big by \eqref{mmp:sections}, so a Mori fiber
space cannot occur: an effective representative of a positive
multiple cannot have negative degree on curves covering general
fibers.  Each stage is consequently a finite, possibly empty,
sequence ending at a nef adjoint.  Standard MMP preservation gives
the next klt pair and a projective $\mathbb Q$-factorial model.
The rational divisor $K_{V_j}+t_jB_{V_j}$ is nef and big, hence
semiample by the klt base point free theorem
\cite[Theorem~3.3]{KM98}.

\smallskip\noindent
\emph{Canonical steps.}
To apply Lemma~\ref{mmp:stabilization}, we must identify every
adjoint step as a canonical MMP step.  This requires canonical
negativity on the contracted ray and, for a flip, canonical
ampleness on the flipped side.
On its source $U$, let $D=K_U+t_{j+1}B_U$.  The contracted ray
$R_0$ satisfies
\[
 D\cdot R_0<0,\qquad
 (D+\lambda C)\cdot R_0=0,\qquad \lambda>0,
\]
where $C=(t_j-t_{j+1})B_U$ is the current scaling divisor.
Therefore $B_U\cdot R_0>0$ and $K_U\cdot R_0<0$.
This is the usual ray calculation for a directed program
\cite[proof of Corollary~1.3.3]{BCHM10}, and it already identifies
each divisorial step as a $K$-negative divisorial contraction.

For a small contraction $h:U\to R$, the adjoint flip makes $D^+$
relatively ample.  We show that $K_{U^+}$ is relatively ample as
well.  Choose the positive rational
number
\[
 b=\frac{K_U\cdot R_0}{D\cdot R_0}.
\]
The divisor $K_U-bD$ is numerically trivial over $R$ and descends
up to rational linear equivalence.  To see this, a Cartier
multiple is relatively semiample by the relative klt base point
free theorem for the pair $(U,t_{j+1}B_U)$
\cite[Theorem~3.24]{KM98}: subtracting its adjoint $D$ gives
relative ampleness, because $-D$ is relatively ample.  Its
semiample morphism is constant on each connected projective fiber,
as its degree on every fiber curve is zero.  The resulting image
is proper and quasi-finite over $R$, because each connected fiber
maps to one point. It is therefore finite and birational over the
normal $R$, hence equals $R$.
Thus $K_U-bD\sim_{\mathbb Q}h^*A$ for a rational Cartier divisor
$A$ on $R$.

Both morphisms in the adjoint flip are small, so this relation
transforms to
\[
 K_{U^+}-bD^+\sim_{\mathbb Q}(h^+)^*A.
\]
Since $D^+$ is relatively ample and $b>0$, so is $K_{U^+}$.
The adjoint flip is therefore a $K$-flip.  Thus all steps of the
finite stages satisfy the hypotheses of
Lemma~\ref{mmp:stabilization}.

\smallskip\noindent
\emph{The fixed model.}
Concatenate the finite stages and apply
Lemma~\ref{mmp:stabilization}.  Choose an endpoint $Y=V_s$ after
the resulting finite prefix.  If only finitely many nonempty stages
occur, choose $s$ after the last one.  For every $j>s$, the finite
composition $Y\dashrightarrow V_j$ is an isomorphism outside
codimension-at-least-three closed subsets on both models.
This property is preserved under finite composition: on a common
open, remove the next bad set and take its closure on the preceding
model; its dimension does not increase.  Let $Z_j$ be the resulting
closed subset of $Y$.

The divisor $M_j=K_Y+t_jB_Y$ agrees on this common open with
$K_{V_j}+t_jB_{V_j}$.  Choose a multiple Cartier on both varieties
and base point free on $V_j$.  A line bundle on a normal variety
has the same global sections after deleting a subset of codimension
at least two, since it is reflexive.  The two spaces of sections
are therefore identified, and the complete system on $Y$ is
generated outside $Z_j$.  This proves \textup{(ii)} for each $j$.

Finally, $Y$ is projective, $\mathbb Q$-factorial and terminal by
the preceding construction, and a finite composition of these
steps extracts no divisors.  On a common smooth projective
resolution of the finite chain, the effective differences in
Lemma~\ref{mmp:strict} telescope to \eqref{mmp:comparison}.
A divisor mapping onto a prime divisor of $Y$ also corresponds to
a prime divisor of $X$, so its coefficient in $E$ is zero.
Thus $E$ is $q$-exceptional.  For an original prime divisor
contracted on $Y$, terminality gives log discrepancy greater than
one on $Y$, compared with one on $X$.  Its strict transform has
positive coefficient in $E$, proving \textup{(iii)}.
\end{proof}

\section{Two surface tests for nefness}\label{num:section}

Our goal is a numerical criterion for nefness. The first surface will
convert positive intersection square into section growth with one fixed
twist. The second surface will show that a negative curve forces that
square to be positive. Together these statements let the numerical-dimension
hypothesis, imposed later on the original variety, rule out every negative
curve. The arguments apply to rational divisors on an
\(n\)-dimensional variety, for every \(n\geq 3\).  We use the following
precise approximation hypothesis:
\begin{equation}\label{num:approximation}
\begin{gathered}
Y\text{ is a normal projective }n\text{-fold over }\mathbb C,
\quad n\geq3,
\qquad \operatorname{codim}_Y\operatorname{Sing}Y\geq 3,\\
L,B_Y\text{ are }\mathbb Q\text{-Cartier},\qquad
t_j\in\mathbb Q_{>0},\quad t_j\longrightarrow 0,
\qquad M_j=L+t_jB_Y,\\
\text{for each }j\text{ there are a closed }Z_j\subset Y,
\quad\operatorname{codim}_Y Z_j\geq 3,
\quad\text{and }d_j\in\mathbb Z_{>0}\text{ such that}\\
d_jM_j\text{ is Cartier and }|d_jM_j|
\text{ is generated by global sections on }Y\setminus Z_j.
\end{gathered}
\end{equation}
Every positive multiple of \(d_j\) has the same generation property.
The integers \(d_j\) may depend on \(j\).
Restrictions of rational Cartier divisors below mean restrictions as
rational line bundles.  Neither effectiveness nor positivity of \(B_Y\) is part
of \eqref{num:approximation}.

\subsection{A fixed twist detects positive surface square}

\begin{lemma}[Surface growth with a fixed twist]\label{num:growth}
Assume \eqref{num:approximation}, and let \(H\) be a very ample Cartier
divisor on \(Y\).  Fix a projective resolution
\(\pi:\widehat Y\to Y\) which is an isomorphism over the smooth locus.
There are a smooth complete intersection surface
\(S\subset Y\), cut out by \(n-2\) members of \(|H|\), and a Cartier
divisor \(P\) on \(\widehat Y\) with the following properties.
Identify \(S\) with its inverse image in \(\widehat Y\).  Then
\(L|_S\) is nef, so
\begin{equation}\label{num:nonnegative-square}
L^2\cdot H^{n-2}\geq 0.
\end{equation}
If \(dL\) is Cartier, then for every positive integer \(m\) divisible
by \(d\) the restriction map
\begin{equation}\label{num:restriction-surjective}
H^0(\widehat Y,m\pi^*L+P)
\longrightarrow H^0(S,mL|_S+P|_S)
\end{equation}
is surjective.  In particular, as \(m\to\infty\) through multiples of
this one fixed integer \(d\),
\begin{equation}\label{num:quadratic-growth}
h^0(\widehat Y,m\pi^*L+P)
\geq \frac{m^2}{2}\,L^2\cdot H^{n-2}+O(m),
\end{equation}
where the error term depends only on the fixed surface and its divisors.
\end{lemma}

\begin{proof}
Put \(r=n-2\).  Very general members of \(|H|\) have a smooth
complete intersection \(S\) avoiding \(\operatorname{Sing}Y\) and
every \(Z_j\).  Indeed, each of these countably many closed sets has
dimension at most \(n-3\), so the condition that \(r\) hyperplanes
meet any one of them is a proper closed incidence condition.
Over \(\mathbb C\) these can be avoided simultaneously, together
with the closed conditions excluded by Bertini's theorem.
The inverse image of \(S\) is therefore exactly the smooth
complete intersection cut out by the \(r\) pulled-back equations;
there are no additional components over the singular locus.
Each \(M_j|_S\) is semiample, hence nef.  Passing to the limit on
every curve on \(S\) shows that \(L|_S\) is nef and proves
\eqref{num:nonnegative-square}.

Write \(A=\pi^*H\), choose an ample Cartier divisor \(A_0\) on
\(\widehat Y\), and make the fixed choice
\begin{equation}\label{num:fixed-twist}
P=K_{\widehat Y}+rA+A_0.
\end{equation}
Since \(A\) is nef and ampleness is open
\cite[Corollary~1.4.10 and Theorem~1.4.23]{Laz04I}, there exists \(\eta>0\)
such that
\begin{equation}\label{num:ample-errors}
A_0+(r-i)A-\epsilon\pi^*B_Y
\quad\text{is ample whenever }0\leq i\leq r
\text{ and }0\leq\epsilon<\eta.
\end{equation}
For a given positive \(m\) with \(d\mid m\), first choose
\(j=j(m)\) so that \(mt_j<\eta\), and then choose a multiple
\(k\) of \(d_j\) with \(k>m\).  The pullback of
\(|kM_j|\) is generated near \(S\).  A general member \(D_m\)
is smooth there, by Bertini, and
\[
R_m=\frac{m}{k}D_m\geq 0,
\qquad R_m\sim_{\mathbb Q}m\pi^*M_j,
\qquad \mathcal J(R_m)=\mathcal O_{\widehat Y}
\quad\text{near }S.
\]
Here \(\mathcal J(R_m)\) is the multiplier ideal; its asserted
local triviality follows from \(m/k<1\) and the smoothness of
\(D_m\) there.

The divisors \(m\pi^*L+P-iA\) are integral Cartier divisors, and
\[
m\pi^*L+P-iA-K_{\widehat Y}-R_m
\sim_{\mathbb Q}
A_0+(r-i)A-mt_j\pi^*B_Y
\]
is ample for \(0\leq i\leq r\).  Nadel vanishing
\cite[Theorem~9.4.8]{Laz04II}, equivalently
\cite[Theorem~2.4]{LazMI}, gives
\begin{equation}\label{num:nadel}
H^q\!\left(\widehat Y,
\mathcal O_{\widehat Y}(m\pi^*L+P-iA)
\otimes\mathcal J(R_m)\right)=0
\qquad(q>0,\ 0\leq i\leq r).
\end{equation}

Set \(\mathcal F=
\mathcal O_{\widehat Y}(m\pi^*L+P)\otimes\mathcal J(R_m)\).
Tensoring the Koszul complex of the \(r\) equations with
\(\mathcal F\) yields the exact complex
\[
0\longrightarrow\mathcal F(-rA)
\longrightarrow\cdots\longrightarrow
\mathcal F(-A)^{\oplus r}\longrightarrow\mathcal F
\longrightarrow\mathcal O_S(mL|_S+P|_S)
\longrightarrow 0.
\]
At points of \(S\), the equations form a regular
sequence and \(\mathcal F\) is locally free.  Outside \(S\), one
equation is a unit, so the Koszul complex is contractible even after
tensoring with an arbitrary sheaf.  This proves exactness everywhere.
Splitting the complex into short exact sequences and using
\eqref{num:nadel} gives surjectivity on global sections onto its last
term.  The inclusion
\(\mathcal F\subset\mathcal O_{\widehat Y}(m\pi^*L+P)\)
then proves \eqref{num:restriction-surjective}.

Adjunction and \eqref{num:fixed-twist} give
\(P|_S-K_S=A_0|_S\).  Hence \(mL|_S+P|_S-K_S\) is ample.
Kodaira vanishing, the zero-boundary ample case of
\cite[Theorem~2.4]{LazMI}, and surface Riemann--Roch
\cite[Chapter~V, Section~1]{Hartshorne77} give
\[
\begin{split}
h^0(S,mL|_S+P|_S)
={}&\frac{m^2}{2}(L|_S)^2
+\frac m2 L|_S\cdot(2P|_S-K_S)\\
&+\frac12 P|_S\cdot(P|_S-K_S)
+\chi(S,\mathcal O_S).
\end{split}
\]
If \(S\) is disconnected, the formula is read componentwise and
summed.  All its coefficients are fixed before \(m,j,k,D_m\) vary.
Together with \eqref{num:restriction-surjective}, this proves
\eqref{num:quadratic-growth}.
\end{proof}

\subsection{Negative degree forces ambient vanishing}

Lemma~\ref{num:growth} supplies the first surface test. To obtain
nefness once its square is zero, we must rule out curves missed by that
very general surface, including curves in the singular locus. The next
two lemmas will measure their effect on a second surface.

The next elementary estimate retains the cotangent bundle of the
ambient variety.  This is what permits a curve singular at the
point where vanishing will be measured.

\begin{lemma}[Uniform ambient multiplicity]\label{num:jets}
Let \(V\) be a smooth projective complex variety, let
\(\Gamma\subset V\) be an integral curve, and let
\(\iota:N\to V\) be the map from its normalization.
For a section \(s\), write \(\operatorname{ord}_z(s)\) for its order
in the regular local ring of \(V\) at \(z\).
There is a positive constant \(g\), depending only on this map,
such that for every line bundle \(\mathcal D\) on \(V\), every
nonzero section \(s\in H^0(V,\mathcal D)\), and every closed
point \(z\in\Gamma\),
\begin{equation}\label{num:jet-bound}
\operatorname{ord}_z(s)\geq
-\frac{\deg_N\iota^*\mathcal D}{g}.
\end{equation}
\end{lemma}

\begin{proof}
Choose a line bundle \(G\) of positive degree on the fixed smooth
curve \(N\) such that
\(\iota^*\Omega^1_V\) embeds as a subbundle of
\(G^{\oplus a}\) for some \(a\).  Such a choice follows by
global generation of
\((\iota^*\Omega^1_V)^*\otimes G\) for a sufficiently ample
\(G\), followed by dualization.  Set \(g=\deg G>0\).

For a nonzero \(s\), the integer
\[
b=\min_{p\in\Gamma\text{ closed}}\operatorname{ord}_p(s)
\]
is finite and attained.  The leading ambient jet gives a nonzero
section of
\begin{equation}\label{num:leading-jet}
\operatorname{Sym}^b(\iota^*\Omega^1_V)
\otimes\iota^*\mathcal D.
\end{equation}
Here is a justification valid also at singular points of \(\Gamma\).
For \(b\geq1\), the bundles of principal parts on the smooth
variety \(V\) have the locally split exact sequence
\[
0\longrightarrow\operatorname{Sym}^b\Omega^1_V\otimes\mathcal D
\longrightarrow\mathcal P^b(\mathcal D)
\longrightarrow\mathcal P^{b-1}(\mathcal D)
\longrightarrow 0.
\]
The lower jet of \(s\) has zero value at every closed point of
the reduced curve \(\Gamma\). A section of a vector bundle on a
reduced variety with this property is zero, so that lower jet restricts
to the zero section. The principal-parts sequence is locally split
as a sequence of \(\mathcal O_V\)-modules, and therefore remains
exact after restriction to \(\Gamma\) and pullback to \(N\).
The restricted order-\(b\) jet consequently belongs to the left-hand
bundle. Its value is nonzero at a point where the minimum \(b\)
is attained; the induced map on its fiber at any point of \(N\)
above that point is an isomorphism. Its pullback is thus regular and
nonzero.  For \(b=0\), use \(s|_\Gamma\)
directly.  This proves \eqref{num:leading-jet} without assuming
smoothness of \(\Gamma\).

Taking symmetric powers of the subbundle inclusion above embeds
\eqref{num:leading-jet} into a direct sum of copies of
\(G^{\otimes b}\otimes\iota^*\mathcal D\).  A nonzero component
is a nonzero section of this line bundle on \(N\), so
\[
0\leq bg+\deg_N\iota^*\mathcal D.
\]
Since \(\operatorname{ord}_z(s)\geq b\) for every \(z\in\Gamma\),
\eqref{num:jet-bound} follows.
\end{proof}

\subsection{A fixed exceptional curve forces a positive square}

\begin{lemma}[Surface fixed-part estimate]\label{num:fixed-square}
Let \(T\) be an integral projective surface, let
\(f:T_1\to T\) be a projective birational morphism from a smooth
surface, and fix a curve \(e\subset T_1\) contracted by \(f\).
For each \(j\), let \(D_j\) be a rational Cartier divisor on
\(T\), and let \(k_j>0\) be an integer such that \(k_jD_j\)
is Cartier.  Suppose a nonzero linear system of
\(k_jf^*D_j\) is generated outside the inverse image of a finite
subset of \(T\).  Let \(F_j\) be its fixed divisor.
If, for one fixed \(c>0\),
\[
\operatorname{coeff}_e F_j\geq ck_j
\quad\text{for all }j,
\]
then there exists \(\epsilon>0\), independent of \(j\), such that
\begin{equation}\label{num:surface-positive-square}
D_j^2\cdot[T]\geq\epsilon
\quad\text{for all }j.
\end{equation}
No normality assumption on \(T\) is needed.
\end{lemma}

\begin{proof}
Every curve in \(\operatorname{Supp}F_j\) maps to a point of
\(T\), by the generation assumption.  Fix a very ample Cartier
divisor \(H_T\) on \(T\) and an ample Cartier divisor \(J\) on
\(T_1\), and put \(Q=f^*H_T\).  The projection formula gives
\begin{equation}\label{num:orthogonality}
Q^2=H_T^2\cdot[T]>0,
\qquad Q\cdot F_j=f^*D_j\cdot F_j=0.
\end{equation}
The Hodge index theorem on the fixed smooth surface \(T_1\)
\cite[Chapter~V, Section~1]{Hartshorne77} makes the negative intersection
form positive definite on
\(Q^\perp\subset N^1(T_1)_{\mathbb R}\).  Consequently there
is one constant \(K>0\) such that
\[
(J\cdot F)^2\leq K(-F^2)
\qquad\text{for every }F\in Q^\perp.
\]
Effectiveness of \(F_j\) gives
\(J\cdot F_j\geq ck_j(J\cdot e)\).  Thus, with the fixed
positive number \(\epsilon=c^2(J\cdot e)^2/K\),
\begin{equation}\label{num:fixed-negative-square}
F_j^2\leq-\epsilon k_j^2.
\end{equation}

After removal of \(F_j\), the residual linear system has no
fixed curve.  Its square is nonnegative: choose two members
with no common curve and take their effective intersection,
or use a nowhere vanishing member if the system is trivial.
By \eqref{num:orthogonality} and the projection formula,
\[
0\leq(k_jf^*D_j-F_j)^2
=k_j^2(D_j^2\cdot[T])+F_j^2.
\]
Now \eqref{num:fixed-negative-square} proves
\eqref{num:surface-positive-square}.
\end{proof}

\subsection{The nefness criterion}

We now combine the last two estimates. A negative curve forces all
sections to vanish linearly at one fixed ambient point. Restriction to
a surface through that point, followed by one blowup, produces the fixed
exceptional curve required by Lemma~\ref{num:fixed-square}. The choices
are made before the perturbation index varies.

\begin{proposition}[Vanishing surface square implies nefness]
\label{num:nef}
Assume \eqref{num:approximation}.  If, for one very ample Cartier
divisor \(H\) on \(Y\),
\begin{equation}\label{num:zero-square}
L^2\cdot H^{n-2}=0,
\end{equation}
then \(L\) is nef.
\end{proposition}

\begin{proof}
Suppose that an integral curve \(C_0\subset Y\) satisfies
\(L\cdot C_0<0\).  Fix a projective resolution
\(\pi:\widehat Y\to Y\) which is an isomorphism over the smooth
locus.  Choose a closed point \(x\in C_0\) outside the images
of those irreducible components of \(\pi^{-1}(C_0)\) which do
not dominate \(C_0\).  There are only finitely many such images,
and each is a point.  Consequently every point of \(\pi^{-1}(x)\)
lies on a component of \(\pi^{-1}(C_0)\) dominating \(C_0\).

\smallskip\noindent
\emph{A surface through \(x\).}
Choose \(n-2\) very general members
\(H_1,\ldots,H_{n-2}\in|H|\), all passing through \(x\).
They meet properly, and their intersection meets each \(Z_j\)
and \(\operatorname{Sing}Y\) in at most finitely many points.
Indeed, the hyperplanes through \(x\) have no other base point;
successive general cuts avoid all positive-dimensional components
of the relevant intersections.  The countably many conditions
can again be imposed simultaneously over \(\mathbb C\).
Write their intersection cycle as
\[
H_1\cdots H_{n-2}\cdot[Y]=\sum_{i=1}^a a_i[T^{(i)}],
\qquad a_i>0,
\]
where each \(T^{(i)}\) is an integral surface.
The dimension theorem guarantees a component through \(x\);
fix one and denote it by \(T\).

On every \(T^{(i)}\), the system induced by \(|d_jM_j|\) has at
most a finite base locus.  Its square is nonnegative, including
when \(T^{(i)}\) is nonnormal.  To check this directly, choose a
section nonzero at the generic point, and then choose a second
section not vanishing identically on any curve of the first
zero divisor.  The resulting Cartier intersections with
\([T^{(i)}]\) form an effective zero-cycle.  If the first section
has no zero, the restricted line bundle is trivial and its
square is zero.  These are the usual Cartier intersection and
projection formulas on cycles \cite[Chapter~2]{Fulton98}.
It follows that
\[
M_j^2\cdot[T^{(i)}]\geq0,
\qquad
\sum_{i=1}^a a_i(M_j^2\cdot[T^{(i)}])
=M_j^2\cdot H^{n-2}\longrightarrow 0
\]
by \eqref{num:zero-square}.  In particular,
\begin{equation}\label{num:special-surface-limit}
M_j^2\cdot[T]\longrightarrow 0.
\end{equation}

\smallskip\noindent
\emph{Fixed geometry above \(x\).}
The surface \(T\) meets \(\operatorname{Sing}Y\) only in finitely
many points.  Its strict transform \(T_{\widehat Y}\) is thus
birational to \(T\) and maps onto it.  Choose any
\(z\in T_{\widehat Y}\) above \(x\).
There is an integral curve \(\Gamma\subset\widehat Y\)
through \(z\) dominating \(C_0\).  To see this, take an
irreducible component \(A\subset\pi^{-1}(C_0)\) through \(z\).
By the choice of \(x\), it dominates \(C_0\).  If
\(\dim A=d\), its fiber \(A_x\) is a proper closed subset
of dimension at most \(d-1\).  General very ample cuts through
\(z\), in number \(d-1\), cut \(A_x\) to dimension at most
zero and leave a curve component through \(z\) on \(A\).
That component cannot lie in \(A_x\), so it dominates \(C_0\).
For \(d=1\), take \(\Gamma=A\).

Let \(\iota:N\to\widehat Y\) denote the normalization map of
\(\Gamma\), followed by inclusion.
Resolve \(T_{\widehat Y}\) by a smooth projective surface
\(T_0\to T_{\widehat Y}\), choose a point \(w\in T_0\)
above \(z\), and blow it up.  Write \(T_1\) for the resulting
surface, \(e\subset T_1\) for the exceptional curve of this
blowup, and \(f:T_1\to T\) for the composite map:
\[
\begin{tikzcd}[column sep=large,row sep=large]
 T_1 \arrow[r,"\mathrm{Bl}_w"] \arrow[d,"f"'] &
 T_0 \arrow[r] &
 \widehat Y \arrow[d,"\pi"] &
 N \arrow[l,"\iota"'] \arrow[d] \\
 T \arrow[rr,hook] & & Y & C_0 \arrow[l,hook]
\end{tikzcd}
\]
The images of \(T_0\) and \(N\) in \(\widehat Y\) meet at
\(z\).  The exceptional curve \(e\), all these varieties, and
all these maps are fixed throughout the remaining argument.

\smallskip\noindent
\emph{The contradiction on the fixed surface.}
The projection formula and \(M_j\to L\) give, for one fixed
\(\delta>0\),
\begin{equation}\label{num:negative-degree}
\deg_N\iota^*\pi^*M_j\leq-\delta
\qquad\text{for all sufficiently large }j.
\end{equation}
In fact this degree equals the positive degree of
\(N\to C_0\) times \(M_j\cdot C_0\).
For each such \(j\), choose any positive multiple \(k_j\) of
\(d_j\).  Lemma~\ref{num:jets}, applied on the fixed
\(\widehat Y\) and \(\Gamma\), shows that every nonzero
pullback section
\(s'\in H^0(\widehat Y,k_j\pi^*M_j)\) of a section on \(Y\)
satisfies
\begin{equation}\label{num:uniform-point-order}
\operatorname{ord}_z(s')\geq ck_j,
\qquad c=\delta/g>0.
\end{equation}
The constant \(c\) depends only on the fixed normalized curve and
the negative-degree margin in \eqref{num:negative-degree}.

The sections on \(Y\) induce a nonzero linear system of
\(k_jf^*(M_j|_T)\), generated outside
\(f^{-1}(T\cap Z_j)\).  Nonzeroness follows from generation
at points of \(T\setminus Z_j\).

Every nonzero section in this induced system has order at least
\(ck_j\) at \(w\) before the last blowup.  Indeed, the local
ring map
\(\mathcal O_{\widehat Y,z}\to\mathcal O_{T_0,w}\)
sends \(\mathfrak m_z^b\) into \(\mathfrak m_w^b\);
apply \eqref{num:uniform-point-order} to a section inducing it.
Sections whose restriction is identically zero do not contribute
to the induced system.  Thus its fixed divisor \(F_j\) satisfies
\(\operatorname{coeff}_eF_j\geq ck_j\).
All the hypotheses of Lemma~\ref{num:fixed-square} hold with
\(D_j=M_j|_T\).  Lemma~\ref{num:fixed-square} gives a fixed \(\epsilon>0\) such
that
\[
M_j^2\cdot[T]\geq\epsilon
\qquad\text{for all sufficiently large }j,
\]
contrary to \eqref{num:special-surface-limit}.  Therefore no
curve \(C_0\) of negative \(L\)-degree exists, and \(L\) is nef.
\end{proof}

\section{The canonical divisor}
\label{main:proof-section}

We now apply the preceding results to the canonical divisor. The passage back to the original smooth variety is recorded explicitly, because the numerical-dimension hypothesis concerns a fixed ample twist on \(X\).

\begin{proof}[Proof of Theorem~\ref{main:theorem}]
Apply Proposition~\ref{mmp:approximation} to obtain a finite canonical MMP
\[
X\dashrightarrow Y
\]
and the divisors \(L=K_Y\), \(B_Y\), and \(M_j=L+t_jB_Y\), with generated multiples outside closed subsets \(Z_j\) of codimension at least three. The variety \(Y\) is terminal and therefore smooth in codimension two. These are the standing hypotheses of the numerical results.

Fix a very ample Cartier divisor \(H\) on \(Y\) and a projective resolution \(\pi:\widehat Y\to Y\) that is an isomorphism over the smooth locus. Lemma~\ref{num:growth} gives a fixed Cartier divisor \(P\) on \(\widehat Y\). Choosing one positive integer \(r_0\) for which \(r_0L\) is Cartier, we obtain
\begin{equation}\label{main:quadratic}
h^0(\widehat Y,m\pi^*L+P)
\ \ge\ \frac{m^2}{2}L^2\cdot H^{n-2}+O(m)
\qquad (r_0\mid m).
\end{equation}
It also gives \(L^2\cdot H^{n-2}\ge0\).

Suppose that this intersection number is positive. Choose a smooth projective common resolution \(W\) dominating \(\widehat Y\) and all the models in the finite MMP prefix, with maps
\[
r:W\longrightarrow\widehat Y,\qquad u:W\longrightarrow X.
\]
Proposition~\ref{mmp:approximation} gives
\[
u^*K_X-r^*\pi^*L\ge0.
\]
For every positive multiple \(m\) of \(r_0\), this effective difference gives an injection
\[
H^0(\widehat Y,m\pi^*L+P)
\hookrightarrow H^0(W,mu^*K_X+r^*P).
\]

Set \(P_X=u_*r^*P\). This is a fixed integral Weil divisor, hence Cartier on the smooth \(X\). A rational section \(f\) in the right-hand space satisfies
\[
\operatorname{div}_W(f)+mu^*K_X+r^*P\ge0.
\]
Pushing this inequality forward gives
\[
\operatorname{div}_X(f)+mK_X+P_X\ge0.
\]
Thus there is an injective linear map
\[
H^0(W,mu^*K_X+r^*P)
\hookrightarrow H^0(X,mK_X+P_X).
\]
Choose a single ample Cartier divisor \(A_X\) for which \(A_X-P_X\) has a nonzero section. Multiplication by that fixed section gives
\[
h^0(X,mK_X+A_X)\ge h^0(\widehat Y,m\pi^*L+P).
\]
By \eqref{main:quadratic}, the left-hand side is bounded below by \(cm^2\), for a fixed \(c>0\), along all sufficiently large multiples of \(r_0\). Consequently
\[
\limsup_{m\to\infty}\frac{h^0(X,mK_X+A_X)}{m^2}>0,
\]
contradicting \(\kappa_\sigma(X,K_X)=1\) and the convention \eqref{main:sigma}. It follows that
\[
L^2\cdot H^{n-2}=0.
\]

Proposition~\ref{num:nef} now implies that \(L=K_Y\) is nef, including on curves contained in the singular locus. All other conclusions of Theorem~\ref{main:theorem}, including the effective exceptional comparison \eqref{main:comparison} and its strict support condition, are supplied by the finite MMP in Proposition~\ref{mmp:approximation}.
\end{proof}

\bibliographystyle{amsalpha}
\bibliography{references}
\end{document}
